\section{The argument in outline}\label{sec:overview}

This section sets out the chain of the proof and the part of the paper in which each step is proved, and
describes the objects of the argument and the way the three theorems combine; precise
definitions and statements follow in the later parts, to which we refer. Throughout, the parameters
$q_0,d$ and the data of \cref{def:data} are fixed: scales $x\le x_1\le b$, the number $m\asymp\sqrt x$
of lifts, ranks $r\asymp x^5$, $k=r+1-m$ and $t=5r$, the CM field $E=K(\zeta_{q_1})$ with its group
$\Phi$ of $f$ infinite places, the weights $\kappa(\tau)$, and the primes $p_0,p_1,N$.

\subsection{The chain of the proof}\label{ss:overview-chain}
The proof is a chain of implications, each of which is a theorem proved in a definite part of the
paper. \Cref{fig:chain} shows the chain, and \cref{tab:parts} the parts of the paper and their roles.

\begin{figure}[t]
\centering
\small
\begin{tikzpicture}[x=1mm,y=1mm,
  box/.style={draw,rounded corners=1mm,line width=.4pt,align=flush center,inner sep=1.4mm,text width=#1,
    execute at begin node={\hyphenpenalty=10000\exhyphenpenalty=10000}},
  box/.default=33mm,
  imp/.style={-{Stealth[length=1.7mm,width=1.3mm]},line width=.4pt}]
\node[box=72mm,minimum height=30mm] (main) at (-21.5,0) {Theorem~\ref{thm:main-grh}:\\ the Generalized Riemann Hypothesis\\ for Dirichlet $L$-functions\\
  Theorem~\ref{thm:grh-half-plane}, \cref{sec:main}};
\node[box,minimum height=30mm] (rh) at (43,0) {Theorem~\ref{thm:main-rh}:\\ the Riemann\\ Hypothesis\\ Theorem~\ref{thm:zeros-on-line}, by the\\ functional equation\\ (\cref{ss:B-zero-symmetry})};
\draw[imp] (main.east) -- (rh.west);
\node[box] (crit) at (-43,-31) {One-sided criterion\\ \cref{thm:one-sided-criterion}\\ \S\S\ref{sec:transform}--\ref{sec:criterion}\\ \cref{ss:B-analytic-intro} ff.};
\node[box=38mm] (bound) at (0,-31) {One-sided bound\\ $R_{q_0}(y)\le y^{-1/2+\ve}$\\ \cref{thm:one-sided-bound}};
\node[box] (pnt) at (43,-31) {Prime number\\ theorem for the\\ remainders; the data\\ \S\S\ref{sec:pnt}, \ref{sec:data}};
\node[box=38mm] (ineq) at (0,-55) {Main inequality\\ $m\,\rem\le fa\log^5b$\\ \cref{thm:main-inequality}; \cref{sec:lattices}};
\node[box] (A) at (-43,-81) {Theorem~\ref{thm:A}:\\ short forms\\ \cref{part:upper}\\ \S\S\ref{sec:spaces}--\ref{sec:elementary-divisors}};
\node[box=38mm] (C) at (0,-81) {Theorem~\ref{thm:C}:\\ integral forms are long\\ \cref{part:lower}\\ \S\S\ref{sec:one-lift-norm}--\ref{sec:theorem-C}};
\node[box] (B) at (43,-81) {Theorem~\ref{thm:B}:\\ enough forms\\ \cref{part:spanning}\\ \S\S\ref{sec:moment}--\ref{sec:frames}};
\node[box] (P2) at (-43,-108) {Derivatives of\\ theta functions\\ \cref{part:derivatives}\\ \S\S\ref{sec:field}--\ref{sec:sharpness}};
\node[box=38mm] (P6) at (0,-108) {Congruences and\\ regular twists\\ \cref{part:congruences}\\ \S\S\ref{sec:galois-lifts}--\ref{sec:fit}};
\node[box] (ds) at (43,-108) {Discrete series\\ with a large gap\\ \cref{app:ds}};
\node[box=119mm] (bg) at (0,-130) {Background results in the form used, with their sources: \cref{app:background}\\
  used throughout; the analytic ones, \cref{ss:B-analytic-intro}--\ref{ss:B-zero-symmetry}, are proved in full};
\draw[imp] (crit.north) -- (crit.north|-main.south);
\draw[imp] (bound.north) -- (bound.north|-main.south);
\draw[imp] (pnt.west) -- (bound.east);
\draw[imp] (ineq.north) -- (bound.south);
\draw[imp] (A.north) |- (ineq.west);
\draw[imp] (B.north) |- (ineq.east);
\draw[imp] (C.north) -- (ineq.south);
\draw[imp] (P2.north) -- (A.south);
\draw[imp] ([xshift=11mm]P2.north) -- ++(0,3.5) -| ([xshift=-12mm]C.south);
\draw[imp] (P6.north) -- (C.south);
\draw[imp] ([xshift=-11mm]ds.north) -- ++(0,3.5) -| ([xshift=12mm]C.south);
\draw[imp] (ds.north) -- (B.south);
\end{tikzpicture}
\caption{The chain of the proof. Arrows lead from the results used to the result proved; boxes name
the sections containing the proofs.}\label{fig:chain}
\end{figure}

\begin{enumerate}[label=\textup{(S\arabic*)},leftmargin=2.6em]
\item\label{it:chain-main} \emph{The main theorems.} Theorems~\ref{thm:main-grh} and~\ref{thm:main-rh}
are \cref{thm:grh-half-plane,thm:zeros-on-line}. The second follows from the first: by the functional
equation (\cref{ss:B-zero-symmetry}) a zero $s$ of $L(\cdot,\chi)$ with $0<\re s<1$ gives the zero
$1-s$ of $L(\cdot,\bar\chi)$, both have real part at most $1/2$, and so $\re s=1/2$; $\zeta$ is the
$L$-function of the character modulo $1$.
\item \emph{Nonvanishing from a one-sided bound.} \Cref{thm:grh-half-plane} combines the criterion
\cref{thm:one-sided-criterion}, by which $L(\rho,\psi)\ne0$ for every character $\psi$ modulo $q_0$
and every $\rho\ne1$ with $\re\rho>\beta$ as soon as $R_{q_0}(x)\le x^{\beta-1}$ for large $x$ and for
every prime $d$ in a residue class, with the one-sided bound \cref{thm:one-sided-bound}; passing from
$q_0$ to the modulus of $\chi$ changes $L(s,\chi)$ by Euler factors only, and for $\re s\ge1$ the
nonvanishing is classical. The criterion rests on
Landau's theorem for Dirichlet series with nonnegative coefficients and on the choice of $d$
(\cref{sec:transform,sec:criterion}).
\item \emph{The bound at one scale.} \Cref{thm:one-sided-bound} follows from \cref{prop:R-at-x}, which
bounds $R_{q_0}(x)$ by $\totient(q_0)\log^6b/m$ and the remainders at the second scale $x_1$. These
are small by the prime number theorem in progressions (\cref{thm:pnt-Rq}); data with $m\ge\sqrt x$
exist for every large $x$ (\cref{thm:data-exist}); and \cref{cor:slow-variation} passes from integers
to all real $y$.
\item \emph{The main inequality.} \Cref{prop:R-at-x} is the expression of $R_{q_0}(x)$ through the total
remainder (\cref{prop:total-remainder}(c)) together with the main inequality
$m\,\rem\le fa\log^5b$ (\cref{thm:main-inequality}).
\item \emph{Three theorems.} Theorem~\ref{thm:B} provides a family of short forms spanning almost all
of the target; Theorem~\ref{thm:A} bounds its covolume and its index from above, and
Theorem~\ref{thm:C} bounds the same quantities from below. In the comparison the main amount $\hmain$
cancels, and $m\,\rem$ remains, bounded by the errors (\cref{ss:main-inequality}).
\item \emph{Theorem~\ref{thm:A}.} The seesaw identity (\cref{thm:seesaw}) writes the chain of lifts of a
Poincar\'e image as a composite lift of a source form. The size of the composite lift
(\cref{thm:composite-size}) and its divisibility (\cref{thm:divisibility}), computed from the
derivatives of theta functions of \cref{part:derivatives}, give part~(a); part~(b) is the divisibility
read as an index (\cref{sec:size}).
\item \emph{Theorem~\ref{thm:B}.} By \cref{thm:one-translate,thm:few-translates} it suffices that the
projected Poincar\'e map be a near isometry and that its translates have a small second moment; both
are \cref{thm:spanning-inputs}, which rests on the unfolding of Poincar\'e series into Bessel periods
(\cref{thm:unfolding}), the filling theorem (\cref{thm:filling}) with its main term and remaining terms
(\cref{sec:main-term,sec:remaining}), and the equality of multiplicities (\cref{thm:equal-multiplicities}).
\item \emph{Theorem~\ref{thm:C}.} It is \cref{thm:twist-for-each-prime} applied to the twisted lattice of
the target. The inputs are auxiliary forms (\cref{thm:auxiliary-forms}); the factorisation of the
volume of the chain into the volumes of single lifts and ratios of Hecke values
(\cref{thm:factorisation}); and the height of each single-lift ratio (\cref{thm:cost-one-lift}) for a
regular family of twists (\cref{thm:regular-family}), which exists by the congruences of
\cref{part:congruences} (\cref{thm:denominator-bound,thm:fit-one-prime}).
\end{enumerate}

\emph{Closure.} Every theorem, proposition, lemma and corollary of the paper is used in the proof of
Theorem~\ref{thm:main-grh} or Theorem~\ref{thm:main-rh}: it is reached from their proofs by following
the references made in proofs and statements. Each of them is proved in the paper, except the
background results, which are stated with their sources: those of \cref{app:background}, the
analytic ones among them being proved in full there, and the facts on CM curves collected in
\cref{thm:cm-background}. The references between proofs, followed item by item, form no cycle, so the
statements can be proved in an order in which each proof uses only statements already proved. In
\texttt{grh-lean} each statement corresponds to declarations listed, section by section, in
\cref{app:lean}, among them declarations used by the five theorems of \cref{ss:intro-theorems}.

\begin{table}[t]
\centering
\small
\caption{The parts of the paper.}\label{tab:parts}
\begin{tabular*}{\textwidth}{@{}l@{\extracolsep{\fill}}>{\raggedright\arraybackslash}p{.35\textwidth}>{\raggedright\arraybackslash}p{.38\textwidth}@{}}
\toprule
part & results & used for\\
\midrule
\Cref{part:reduction}, \S\S\ref{sec:transform}--\ref{sec:main} & criterion, prime number theorem for the remainders, twisted lattices, data, main inequality & Theorems~\ref{thm:main-grh} and~\ref{thm:main-rh}\\
\Cref{part:derivatives}, \S\S\ref{sec:field}--\ref{sec:sharpness} & divisibility of derivatives of theta functions of CM curves, and its sharpness & sizes and divisibility in Theorems~\ref{thm:A} and~\ref{thm:C}\\
\Cref{part:upper}, \S\S\ref{sec:spaces}--\ref{sec:elementary-divisors} & spaces and chains, seesaw identity, divisibility and size of the composite lift & Theorem~\ref{thm:A}; the spaces in Theorems~\ref{thm:B} and~\ref{thm:C}\\
\Cref{part:spanning}, \S\S\ref{sec:moment}--\ref{sec:frames} & equal multiplicities, unfolding, filling theorem, main and remaining terms & Theorem~\ref{thm:B}\\
\Cref{part:lower}, \S\S\ref{sec:one-lift-norm}--\ref{sec:theorem-C} & Rallis' formula, exchange identity, norm of the chain, auxiliary forms, cost of one lift & Theorem~\ref{thm:C}\\
\Cref{part:congruences}, \S\S\ref{sec:galois-lifts}--\ref{sec:fit} & Galois representations of lifts, Ribet's method, denominators, Selmer data, twisting & regular families of twists in Theorem~\ref{thm:C}\\
\Cref{app:ds} & discrete series with a large gap & \cref{part:spanning,part:lower}\\
\Cref{app:background} & background results with their sources & throughout\\
\bottomrule
\end{tabular*}
\end{table}

\subsection{Spaces and forms}\label{ss:overview-spaces}
Let $E^+$ be the maximal totally real subfield of $E$. All hermitian spaces are over $E$ and
\emph{definite}: their unitary groups are compact at every infinite place of $E^+$. Automorphic forms on
such a group are \emph{algebraic modular forms} in the sense of Gross: functions on a finite set of
double cosets with values in a finite-dimensional representation, the weight. They carry a natural
integral structure at every prime, and their spaces carry Petersson inner products given by finite
sums.

\begin{itemize}[leftmargin=1.5em]
\item The spaces $X_0,X_1,\dots,X_m$ have ranks $n_j=k+j(t+1)$. The \emph{source} $\mathcal S$ is a
space of forms on $\U(X_0)$ of a weight concentrated at the $a$-places, ramified at the level $N$.
\item The spaces $Y_0,Y_1,\dots,Y_m$ have ranks $r+jt$, with $Y_m\subset X_m$ of corank one. The
\emph{target} $\mathcal T$ is the space of \emph{good} forms on $\U(Y_0)$: forms of principal level at
the chosen places above $p_0,p_1$ cut out by an idempotent $e_0$ of the group algebra of
$\bar G=\prod_w\GL_r(\F_w)$, built from Green's characters of finite general linear groups.
\item For $1\le j\le m$, the pair $(W_j,V_j)$ of spaces of ranks $r$ and $r+t$ is a \emph{descent pair}:
$W_j$ differs from $Y_0$ only in its signs at the real places, and $V_j$ is definite.
\end{itemize}

\subsection{Maps}\label{ss:overview-maps}
\begin{itemize}[leftmargin=1.5em]
\item The \emph{Poincar\'e map} $\mathcal P\colon\mathcal S\to\{\text{forms on }\U(Y_0)\}$ sends a source
form to the form that represents its Bessel period in the Petersson product. After projection by
$e_0$, translation by $g\in\bar G$ and a normalisation, it gives the maps
$\mathfrak p_g\colon\mathcal S\to\mathcal T$ of Theorems~\ref{thm:A} and~\ref{thm:B}.
\item For each \emph{twist}, an anticyclotomic Hecke character $\xi$ of $E$ of weight $0$ whose tame part
has nontrivial powers $\xi^1,\dots,\xi^m$ at the places of its modulus, the \emph{chain}
$\Theta_\xi\colon\mathcal T\to\{\text{forms on }\U(Y_m)\}$ is the composite of the $m$ theta lifts
$Y_{j-1}\to Y_j$ with the characters $\mu_j^0\xi^j$ of $\U(1)$, of weights $-\kappa_j$.
\item The \emph{composite lift} $B_\xi$ on $\mathcal S$ is a finite combination of chains of theta
lifts $X_0\to X_1\to\cdots\to X_m$, restricted to $Y_m$. The \emph{seesaw identity}
(\cref{part:upper}) states that on ramified source forms
\[
  \Theta_\xi\circ\mathcal P=c_{\rm ss}\,B_\xi
\]
for a constant $c_{\rm ss}\neq0$ independent of $\xi$.
\end{itemize}

The two sides of the seesaw identity are read in different ways. The right side is built from theta
functions attached to the vectors of lattices of $X_j$; its values are combinations of derivatives
(\emph{jets}) of theta functions of CM elliptic curves, whose sizes and divisibilities can be computed.
The left side is a chain of lifts of a target form, whose norms are governed by $L$-values.

\subsection{Theorem A}\label{ss:overview-A}
Each of the $m$ lifts of the composite lift differentiates theta functions: to order $a$ along the
$a$-places, and to order $\kappa(\tau)$ in each of about $k$ further directions, at every place $\tau$.
Two estimates are made for the values of $B_1$ on a source form.
\begin{itemize}[leftmargin=1.5em]
\item \emph{Size.} At every conjugate, a derivative of order $j$ costs at most $\sinf(j)$, with
$\sinf$ as in \eqref{eq:sinf}.
\item \emph{Divisibility.} At each prime $\ell$ the values are divisible by $\ell$ to the power
$g_\ell$, the \emph{gain}, up to an error whose total, weighted by $\log\ell$, is at most the error unit.
The gain is the sum over the lifts and the places of the local parts of the prime sums $\sfin$, so that
$\sum_\ell g_\ell\log\ell$ is computed by prime sums.
\end{itemize}
Dividing $B_1$ by the rational number $\prod_\ell\ell^{g_\ell}$ gives integral forms, and the gains
cancel between the two estimates. In the $k$ further directions the size $\sinf(\kappa(\tau))$ is
repaid only by $\sfin(\kappa(\tau),\tau)$, which leaves the main amount
$\hmain=mk\sum_\tau(\sinf-\sfin)$. Along the $a$-places the size is matched by the expected value
$\smain$ of the prime sums, while the divisibility is given by the actual prime sums
$\sfin(a,\tau)$; their difference, summed over the $a$-places and the $m$ lifts, is $m\,\rem$. The
result is the bound $\exp(\hmain-m\,\rem+fa\log^4b)$ of Theorem~\ref{thm:A}(a). Part~(b), the index
bound, records that the divisibility holds on a submodule of small colength in the integral source
forms. The bookkeeping is carried out in \cref{part:upper}.

\subsection{Theorem B}\label{ss:overview-B}
Theorem~\ref{thm:B} asserts that a few translates $\mathfrak p_g(\mathcal S)$, $g\in \Sigma$, already span
almost all of $\mathcal T$. It follows from a second-moment bound for the operators $\mathcal P^*g\mathcal P$:
the mean over $g\in\bar G$ is close to a multiple of the projection to $\mathcal T$, by an unfolding of
Poincar\'e series into Bessel periods, an evaluation of the main term by Siegel's formula and a sieve
over the places, and a bound for the remaining terms by Schur's test (\cref{part:spanning}).

\subsection{Theorem C}\label{ss:overview-C}
Theorem~\ref{thm:C} is an application of \cref{thm:twist-for-each-prime} to the twisted lattice
$\mathcal T$. Its inputs are:
\begin{enumerate}[label=\textup{(\roman*)}]
\item a factorisation of $Z(1)/Z(\xi)$, the ratio of the Gram determinants of the untwisted and the
twisted chain, into the ratios $Z_j(1)/Z_j(\xi^j)$ for the single lifts $W_j\to V_j$, times a power of a
scalar given by ratios of Hecke $L$-values (\cref{part:lower});
\item for each prime place $\iota$, a twist $\xi_\iota$ allowed at the prime of $\iota$ such that each
ratio $Z_j(1)/Z_j(\xi_\iota^j)$ has height at most $\dimT\bigl(fa\log^4b/m-k\,\hdir\bigr)$. The norm of
a single lift is a Hecke $L$-value by the Rallis inner product formula; the height bound combines the
archimedean size of this value with a bound for its denominators, obtained from congruences between
theta lifts and the smallness of certain anticyclotomic Selmer groups (\cref{part:congruences});
\item auxiliary forms that complete a family of integral forms to a basis at a controlled cost.
\end{enumerate}
Multiplying the $m$ single-lift bounds gives a height of about $-\dimT\,\hmain$ for $Z(1)/Z(\xi_\iota)$,
and \cref{thm:twist-for-each-prime} turns it into the lower bound $c+2\delta\ge\dimT(\hmain-O(fa\log^4b))$.

\subsection{Sizes}\label{ss:overview-sizes}
The exponents are chosen so that every error is smaller than the main amount by a power of $x$, and
so that the main inequality yields the exponent $1/2$. \Cref{prop:sizes} lists the sizes. The
choices interact as follows.
\begin{itemize}[leftmargin=1.5em]
\item The main inequality has the form $m\,\rem\le fa\log^5b$, and $\rem$ is about
$a\,\#\Phib\,R_{q_0}(x)$ with $\#\Phib\ge f/(\totient(q_0)\log b)$. Hence $R_{q_0}(x)\ll
\log^6b/m$, and $m\ge\sqrt x$ with $\log b\le2x^\ve$ gives the exponent $-1/2+6\ve$.
\item The number of lifts cannot exceed about $\sqrt x$: a chain of $m$ lifts adds $m$ blocks of
$t+1$ dimensions, and the gap $\Delta=2r^8$ between the weights of consecutive lifts must dominate
the ranks, while all lift weights stay within $x^{42}$ of $\kappa$ (\cref{ss:data}). The relation
between $m$ and $x$ is the quantitative form of the remark, recalled in \cref{ss:intro-method}, that a
chain of $m$ lifts plays the part of vanishing to order $m$.
\item The second scale $x_1\ge x^{10^6}$ makes the prime sums at the places outside $\Phib$ run up to
$x_1$, where the prime number theorem applies; the third scale $b\ge x_1^{10^6}$ makes the weights
large against all other sizes, so that the errors are measured in the unit $\errunit=fa\log^3b$.
\end{itemize}
